% Options for packages loaded elsewhere
\PassOptionsToPackage{unicode}{hyperref}
\PassOptionsToPackage{hyphens}{url}
\documentclass[
]{article}
\usepackage{xcolor}
\usepackage{amsmath,amssymb}
\setcounter{secnumdepth}{-\maxdimen} % remove section numbering
\usepackage{iftex}
\ifPDFTeX
  \usepackage[T1]{fontenc}
  \usepackage[utf8]{inputenc}
  \usepackage{textcomp} % provide euro and other symbols
\else % if luatex or xetex
  \usepackage{unicode-math} % this also loads fontspec
  \defaultfontfeatures{Scale=MatchLowercase}
  \defaultfontfeatures[\rmfamily]{Ligatures=TeX,Scale=1}
\fi
\usepackage{lmodern}
\ifPDFTeX\else
  % xetex/luatex font selection
\fi
% Use upquote if available, for straight quotes in verbatim environments
\IfFileExists{upquote.sty}{\usepackage{upquote}}{}
\IfFileExists{microtype.sty}{% use microtype if available
  \usepackage[]{microtype}
  \UseMicrotypeSet[protrusion]{basicmath} % disable protrusion for tt fonts
}{}
\makeatletter
\@ifundefined{KOMAClassName}{% if non-KOMA class
  \IfFileExists{parskip.sty}{%
    \usepackage{parskip}
  }{% else
    \setlength{\parindent}{0pt}
    \setlength{\parskip}{6pt plus 2pt minus 1pt}}
}{% if KOMA class
  \KOMAoptions{parskip=half}}
\makeatother
\usepackage{longtable,booktabs,array}
\usepackage{calc} % for calculating minipage widths
% Correct order of tables after \paragraph or \subparagraph
\usepackage{etoolbox}
\makeatletter
\patchcmd\longtable{\par}{\if@noskipsec\mbox{}\fi\par}{}{}
\makeatother
% Allow footnotes in longtable head/foot
\IfFileExists{footnotehyper.sty}{\usepackage{footnotehyper}}{\usepackage{footnote}}
\makesavenoteenv{longtable}
\setlength{\emergencystretch}{3em} % prevent overfull lines
\providecommand{\tightlist}{%
  \setlength{\itemsep}{0pt}\setlength{\parskip}{0pt}}
\usepackage{bookmark}
\IfFileExists{xurl.sty}{\usepackage{xurl}}{} % add URL line breaks if available
\urlstyle{same}
\hypersetup{
  hidelinks,
  pdfcreator={LaTeX via pandoc}}

\author{}
\date{}

\begin{document}

\section{Autopoietic Cohomology: Iterative Obstruction Repair on Causal
Posets}\label{autopoietic-cohomology-iterative-obstruction-repair-on-causal-posets}

\textbf{Author:} Jeremy H. Carroll \textbf{Date:} March 2026
\textbf{Version:} v1.0 (Pre-Print)

\begin{center}\rule{0.5\linewidth}{0.5pt}\end{center}

\subsection{Abstract}\label{abstract}

\textbf{Autopoietic cohomology describes a process in which unresolved
\(\ell^1\) coboundary defects induce the minimal augmentation of the
underlying causal structure required for their resolution.}

We define a discrete dynamical system --- the \textbf{Autopoietic
Iterator} --- that executes this structural extension by iteratively
repairing cohomological obstructions on a causal poset. At each step,
the system detects the leading obstruction class
\([o_n] \in H^n(\mathcal{P}; F)\), computes a minimal-cost transport
extension \(\Lambda_n\) via \(\ell^1\) optimization, and forms the
augmented poset
\(\mathcal{P}_{n+1} = \mathcal{P}_n \coprod_{\partial} \Lambda_n\).
Because defect resolution is achieved through structural extension
rather than internal reconfiguration, the construction is strictly
monotone (the poset only grows) and the coboundary norm is
non-increasing. We observe that the repair LP admits a formulation
closely related to the Kantorovich optimal transport problem, making
each repair step a discrete W₁ transport optimization.

This construction synthesizes the obstruction results of {[}1{]}, the
\(\ell^1\) classification theorems of {[}2, 3{]}, the cohomological
structure of {[}4{]}, and the universal framework of {[}5{]}. We propose
--- as a suggestive analogy --- that this iterative repair process
provides a candidate dynamical mechanism for Deutsch's Constructor
Theory {[}6{]}, in which physical evolution corresponds to sequential
resolution of topological inconsistencies.

\begin{center}\rule{0.5\linewidth}{0.5pt}\end{center}

\subsection{1. Introduction}\label{introduction}

\subsubsection{1.1 Background}\label{background}

In the preceding papers:

\begin{itemize}
\tightlist
\item
  \textbf{Paper 000} {[}1{]} and \textbf{Paper 003} {[}4{]} established
  that projection dynamics generate coboundary defects, persisting with
  canonical cohomological structure.
\item
  \textbf{Paper 001} {[}2{]} and \textbf{Paper 002} {[}3{]} classified
  admissible diagnostics on Banach presheaves and finite graphs, showing
  \(\ell^1\) total variation is uniquely selected.
\item
  \textbf{Paper 004} {[}5{]} synthesized these into the Universal
  Obstruction Theory framework, extending to operator-valued settings
  and W₁ duality.
\end{itemize}

The missing piece is \textbf{dynamics}: given a poset with non-trivial
obstructions, what process resolves them, and what structure does it
generate?

\subsubsection{1.2 The Autopoietic Structural
Principle}\label{the-autopoietic-structural-principle}

We formalize the principle that \textbf{defect resolution is not
internal optimization---it is structural extension}. If an \(\ell^1\)
defect cannot be reduced within an existing poset, geometric consistency
demands the augmentation of the underlying morphism structure.

We define the \textbf{Autopoietic Iterator} as the formal discrete
dynamical system mapping this process: it iteratively detects
unresolvable obstruction classes and augments the causal poset to
resolve them via minimal \(\ell^1\) transport geometry.

This paper defines the iterator precisely and proves its basic
structural properties.

\begin{center}\rule{0.5\linewidth}{0.5pt}\end{center}

\subsection{2. Mathematical
Construction}\label{mathematical-construction}

\subsubsection{2.1 The Autopoietic
Iterator}\label{the-autopoietic-iterator}

Let \(\mathcal{P}_0\) be a finite causal poset with presheaf
\(F_0: \mathcal{P}_0^{\text{op}} \to \mathbf{Ban}\), and suppose
\(I(s) = \sum_e \|\delta_e(s)\|_1 > 0\) (non-trivial coboundary norm).

\textbf{Definition 2.1 (Autopoietic Iterator).} The sequence
\(\{\mathcal{P}_n\}_{n \geq 0}\) is defined recursively:

\textbf{(i) Detect.} Compute the leading obstruction class
\([o_n] \in H^{k_n}(\mathcal{P}_n; F_n)\), defined as the class with the
maximal \(\ell^1\) norm contribution to the global stable norm.

\textbf{(ii) Optimize.} Compute the optimal repair morphism:

\[\Lambda_n = \arg\min_{\Lambda} \; I_{\mathcal{P}_n \cup \Lambda}(s)\]

where the minimization is evaluated over all admissible candidate
additions
\(\Lambda \subset (V \times V) \setminus E_{\text{cov}}(\mathcal{P}_n)\)
satisfying boundary compatibility conditions restricted to
\(\partial(\text{supp}([o_n]))\), and \(I\) is the \(\ell^1\) coboundary
norm.

\textbf{(iii) Augment.} Form the pushout:

\[\mathcal{P}_{n+1} = \mathcal{P}_n \coprod_{\partial(\text{supp}([o_n]))} \Lambda_n\]

where the pushout is taken along boundary inclusions in the category of
finite posets with monotone maps.

\textbf{(iv) Terminate.} If \(I_{\mathcal{P}_{n+1}}(s) = 0\) (all
obstructions resolved), halt. Otherwise, return to (i).

\textbf{Remark.} The optimization step (ii) is an \(\ell^1\) linear
program. For finite posets, it is solvable via standard linear
programming methods (polynomial-time in the size of the discretized
problem). The pushout in step (iii) produces a strictly larger poset.

\subsubsection{2.2 The Complexity Ratchet}\label{the-complexity-ratchet}

\textbf{Definition 2.2 (Monotone Growth Constraint).} The Autopoietic
Iterator satisfies the \textbf{Complexity Ratchet}: the inclusion
\(\mathcal{P}_n \hookrightarrow \mathcal{P}_{n+1}\) is strict, and all
existing covering relations are preserved. No morphisms may be deleted:

\[E_{\text{cov}}(\mathcal{P}_n) \subseteq E_{\text{cov}}(\mathcal{P}_{n+1}), \qquad |\mathcal{P}_{n+1}| > |\mathcal{P}_n| \text{ whenever } I_{\mathcal{P}_n}(s) > 0\]

\textbf{Remark 2.2 (Modeling assumption).} The monotone growth
constraint is a design choice, not a mathematical necessity. We restrict
attention to monotone augmentations to model discrete irreversibility.
However, it is critical to note that \textbf{monotone growth does not
imply monotone simplification of cohomology}. Introducing new causal
edges to repair a target obstruction can inadvertently introduce
secondary obstructions across the expanded topology.

\subsubsection{2.3 Basic Properties}\label{basic-properties}

\textbf{Proposition 2.1 (Monotone Cost Reduction).} \emph{If
\(\Lambda_n\) is chosen to minimize
\(I_{\mathcal{P}_n \cup \Lambda}(s)\), then:}

\[I_{\mathcal{P}_{n+1}}(s) \leq I_{\mathcal{P}_n}(s)\]

\emph{with equality only if no repair morphism reduces the coboundary
norm. If additionally \(\Lambda_n\) trivializes the leading obstruction
class (i.e., \([o_{k_n}]|_{\mathcal{P}_{n+1}} = 0\)), the corresponding
stable-norm summand vanishes from the cost.}

\textbf{Proof.} The identity morphism \(\Lambda = \emptyset\) (adding no
new edges) is always a feasible choice, yielding
\(I_{\mathcal{P}_n \cup \emptyset}(s) = I_{\mathcal{P}_n}(s)\). Since
\(\Lambda_n\) is chosen as the minimizer over all admissible
\(\Lambda\), we directly obtain
\(I_{\mathcal{P}_{n+1}}(s) \leq I_{\mathcal{P}_n}(s)\).

For the stronger claim: if \([o_{k_n}]|_{\mathcal{P}_{n+1}} = 0\), then
by the cost formula (Theorem 2 of {[}4{]}), the \(k_n\)-th summand
\(\|[o_{k_n}]\|_1\) vanishes in the augmented poset. The remaining
summands \(\|[o_k]\|_1\) for \(k \neq k_n\) may change, but the total
cost remains bounded by the minimality property. \(\square\)

\textbf{Open Questions.}

\begin{enumerate}
\def\labelenumi{\arabic{enumi}.}
\tightlist
\item
  \textbf{Convergence.} Does the iterator terminate in finitely many
  steps for all finite posets, or can it cycle indefinitely? Monotone
  norm decrease does not guarantee termination, since the \(\ell^1\)
  cost takes values in a continuous range and could decrease
  asymptotically. For bounded-degree posets with integer-valued stalks,
  finite termination is guaranteed, but the general case is open.
\item
  \textbf{Uniqueness.} Is the optimal \(\Lambda_n\) unique, or does a
  highly symmetric \(\mathcal{P}_n\) accommodate multiple
  cost-equivalent minimal extensions?
\item
  \textbf{Complexity.} What is the formal computational complexity bound
  of the full discrete trace acting as a function of
  \(|\mathcal{P}_0|\)?
\end{enumerate}

\begin{center}\rule{0.5\linewidth}{0.5pt}\end{center}

\subsection{3. Algorithmic
Interpretation}\label{algorithmic-interpretation}

This framework dictates a dual purpose: first, as a theoretical
foundation, and second, as a concrete, executable computing structure.
The abstract iterator definition provides a directly implementable
algorithmic protocol.

\subsubsection{3.1 Discrete Wasserstein
Descent}\label{discrete-wasserstein-descent}

The repair optimization in step (ii) of the iterator assumes the
structural configuration of the \textbf{Kantorovich optimal transport
problem} {[}9{]}. The stable norm computation:

\[\|[o_n]\|_1 = \min_{\omega} \sum_e |\omega(e)| \quad \text{subject to} \quad \omega = o_n + \delta f\]

is the Kantorovich dual of optimal transport: minimizing the \(\ell^1\)
cost to move defect mass subject to a continuous divergence (coboundary)
constraint.

\begin{longtable}[]{@{}ll@{}}
\toprule\noalign{}
This framework & Optimal transport \\
\midrule\noalign{}
\endhead
\bottomrule\noalign{}
\endlastfoot
Defect cocycle \(o_n\) & Mass imbalance \(\mu - \nu\) \\
Cochain correction \(\delta f\) & Transport plan \(\gamma\) \\
\(\ell^1\) coboundary norm & Wasserstein-1 cost \\
LP dual (max-flow) & Kantorovich dual \\
\end{longtable}

This identification arises naturally from the dual structure of the
\(\ell^1\) optimization. A discrete W₁ transport problem operates
iteratively by moving topological defects toward minimal boundaries.

\textbf{Connection to minimum cut.} Since the \(\ell^1\) coboundary norm
scales proportionally to the graph total variation, each repair
optimization maps back cleanly to a minimum cut extraction on the defect
support graph. Inserting the optimal \(\Lambda_n\) matches inserting a
structural bridge across the min-cut boundary separating conflicting
defect regions. Each descent step evaluates computationally within
polynomial time via established algorithms like push-relabel.

\textbf{Why the \(\ell^1\) Geometry.} The \(\ell^1\) norm naturally
yields sparse, localized, LP-solvable repairs within this framework. An
evaluated \(\ell^2\) integration diffuses defects globally throughout
the entire connected lattice structure. The \(\ell^1\) formulation
provides targeted minimum cuts and transport mechanics, facilitating a
constrained, executable descent.

\subsubsection{3.2 The Omega Engine as
Evidence}\label{the-omega-engine-as-evidence}

An illustrative computational implementation of the iterative formalism
reinforces this. Preliminary observations (not formal benchmarks)
suggest the \textbf{Omega Engine} (\texttt{anomalon\_kernel})
instantiates the Autopoietic Iterator in silico:

\begin{itemize}
\tightlist
\item
  \textbf{Detect:} Locate \(\ell^1\)-maximized cohomological defects
  iteratively across the active graph state.
\item
  \textbf{Optimize:} Execute the associated \(\ell^1\) linear objective
  to chart the sparsest minimum transport routing.
\item
  \textbf{Augment:} Compile the proposed modification into structural
  code boundaries mimicking the posetal pushout.
\end{itemize}

If the iterator strictly operates as a transport optimization, the
engine should exhibit specific, computationally verifiable signatures:
1. \textbf{Defect transport:} Cohomological obstruction density will map
deterministically along minimal cost constraints evaluated iteratively.
2. \textbf{Min-cut concentration:} Codebase or topological graph
adaptations should concentrate overwhelmingly across min-cut margins
rather than randomized surface areas. 3. \textbf{Sparse repair:} Active
topological integrations consistently display finite, targeted revisions
in isolated regions.

These measurable programmatic behaviors provide structural evidence
consistent with the model's formulation.

\begin{center}\rule{0.5\linewidth}{0.5pt}\end{center}

\subsection{4. Connection to the Paper
Series}\label{connection-to-the-paper-series}

The Autopoietic Iterator functionally unifies the preceding anomaly
architecture:

\begin{longtable}[]{@{}
  >{\raggedright\arraybackslash}p{(\linewidth - 2\tabcolsep) * \real{0.3684}}
  >{\raggedright\arraybackslash}p{(\linewidth - 2\tabcolsep) * \real{0.6316}}@{}}
\toprule\noalign{}
\begin{minipage}[b]{\linewidth}\raggedright
Paper
\end{minipage} & \begin{minipage}[b]{\linewidth}\raggedright
Result used
\end{minipage} \\
\midrule\noalign{}
\endhead
\bottomrule\noalign{}
\endlastfoot
{[}1{]} / {[}4{]} --- Obstruction \& Cohomology & Obstructions exist and
persist; \(I(s) > 0\) forms the active initial condition. \\
{[}2{]} / {[}3{]} --- \(\ell^1\) Classification & The cost functional
\(I\) is uniquely determined for boundary isolation. \\
{[}5{]} --- Universal Framework & Wasserstein-1 framework identifies
repair optimization as explicit transport tracking. \\
\end{longtable}

\begin{quote}
\emph{The Autopoietic Iterator evaluates functionally as a discrete
dynamical system minimizing the persistent \(\ell^1\) coboundary norm by
integrating cost-optimal topological modifications.}
\end{quote}

\begin{center}\rule{0.5\linewidth}{0.5pt}\end{center}

\subsection{5. Physical Conjectures}\label{physical-conjectures}

The following interpretations are entirely speculative. They are
proposed as suggestive analogies, motivated exclusively by the
structural topology of the framework, and do not present mathematically
established mappings to conventional physical systems.

\subsubsection{5.1 Time as Ordinal
Iteration}\label{time-as-ordinal-iteration}

\textbf{Conjecture 1 (Speculative).} \emph{Physical time correlates
dimensionally with the ordinal progression count \(n\) of the discrete
Autopoietic Iterator. The irreversibility modeled by discrete
thermodynamic entropy translates topologically via the Complexity
Ratchet constraint.}

This remains a primarily philosophical conjecture. Evaluating it against
standard formulations (like Lorentz invariance and relativistic
causality structures) would demand deriving continuous limit sequences
and demonstrating macroscopic time symmetries currently absent from the
construction.

\subsubsection{5.2 Geometry from Obstruction
Density}\label{geometry-from-obstruction-density}

\textbf{Conjecture 2 (Speculative).} \emph{Macroscopic geometry
evaluates fundamentally from the concentration density determining
continuous obstruction integrations. Deep recursion within a specific
dense obstruction boundary implies severe geometric curvature
anomalies---a topological analogue to gravity.}

Providing this hypothesis formalized rigor necessitates modeling
emergent phenomena mirroring Einstein field configurations, deriving
explicit limits from localized Kan extension concentrations.

\subsubsection{5.3 Decoherence as
Truncation}\label{decoherence-as-truncation}

\textbf{Conjecture 3 (Speculative).} \emph{This serves as a suggestive
analogy, not a derived model of quantum measurement: Wavefunction
collapse heuristically maps to the forced truncation of the
\(\infty\)-groupoid occurring wherever localized processing saturation
mathematically exceeds internal structural capacity boundaries dictated
by local Lieb-Robinson bounds. Exceeding dimension thresholds prompts
sudden truncation down to a singular definable 1-topos representation (a
discrete outcome).}

This represents the broadest metaphysical speculation included herein.
Canonical decoherence dictates strict environmental trace mechanics and
explicit reductions of entangled density boundaries. Correlating this
mathematically requires defining analogous environment structures within
the poset and proving the theoretical reduction accurately reproduces
Born probability thresholds.

\begin{center}\rule{0.5\linewidth}{0.5pt}\end{center}

\subsection{6. Discussion}\label{discussion}

This paper formally branches from Papers 1--4 by transitioning the
established topological metrics into an active dynamic processing
definition.

The mathematical formulation provides: - \textbf{Definition 2.1:}
Defines the active iterator via sequential ℓ¹ defect targeting and
posetal pushouts. - \textbf{Definition 2.2:} Codifies irreversible
evolution explicitly via the Complexity Ratchet assumption. -
\textbf{Proposition 2.1:} Confirms global cumulative optimization
guarantees assuming monotonic edge construction.

Additionally, isolating the Algorithmic Interpretation highlights that
this formalism constructs explicit programmatic pipelines (evidenced by
the AutoP anomaly compiler logic), bridging abstract topologies with
testable discrete behaviors like localized minimum cuts.

Evaluating the physical conjectures depends strictly on solving defining
bounds around iterator termination constraints, extracting computable
continuous analogues, and executing wide-scale computational profiling
of localized engine integrations. It strictly presents a mapped
hypothesis outlining discrete defect-driven thermodynamic and
chronological behaviors.

\begin{center}\rule{0.5\linewidth}{0.5pt}\end{center}

\subsection{7. References}\label{references}

\begin{enumerate}
\def\labelenumi{\arabic{enumi}.}
\tightlist
\item
  Carroll, J. H. (2026). \emph{Single Obstruction Theory: Retraction
  Nonlinearity, \(\ell^1\) Rigidity, and Density Scaling.} (Paper 000).
  Zenodo. https://doi.org/10.5281/zenodo.19151803
\item
  Carroll, J. H. (2026). \emph{The Free \(\ell^1\) Seminorm on Banach
  Presheaf Coboundaries.} (Paper 001). Zenodo.
  https://doi.org/10.5281/zenodo.19152115
\item
  Carroll, J. H. (2026). \emph{Coordinate-Wise Additivity and the
  \(\ell^1\) Norm on Finite Graph Cochains.} (Paper 002). Zenodo.
  https://doi.org/10.5281/zenodo.19152599
\item
  Carroll, J. H. (2026). \emph{Hodge Structure and Gauge Equivalence of
  \(\ell^1\) Defect Fields.} (Paper 003). Zenodo.
  https://doi.org/10.5281/zenodo.19152799
\item
  Carroll, J. H. (2026). \emph{Universal Obstruction Theory: The
  \(\ell^1\) Topological Framework.} (Paper 004). Zenodo.
  https://doi.org/10.5281/zenodo.19154775
\item
  Deutsch, D. (2013). Constructor Theory. \emph{Synthese}, 190(18),
  4331--4359.
\item
  Lurie, J. (2009). \emph{Higher Topos Theory.} Princeton University
  Press.
\item
  Lieb, E. H. and Robinson, D. W. (1972). The finite group velocity of
  quantum spin systems. \emph{Comm. Math. Phys.}, 28(3), 251--257.
\item
  Villani, C. (2009). \emph{Optimal Transport: Old and New.} Springer.
\end{enumerate}

\end{document}
